% Libro 4 · Mundo 4 · Argumentación, intención y tono (app: pu-argumentacion)
\mundo{Argumentación, intención y tono}{pu-argumentacion}{Muchos textos del examen quieren convencer. Leerlos bien es
reconocer su esqueleto (una tesis y las razones que la sostienen), su intención (¿informar, persuadir, criticar,
advertir?) y su tono (la actitud que se oye detrás de las palabras). Y, sobre todo, separar los hechos de las opiniones.}{%
\item \textbf{4.1} Tesis y argumentos
\item \textbf{4.2} Intención y tono
\item \textbf{4.3} Hechos y opiniones}

\leccion{pu-ag-1}{Tesis y argumentos}{Juego de armar el árbol de un argumento.}

\tecnica{La técnica: el árbol del argumento}
Un texto argumentativo tiene una \clave{tesis} (la idea que defiende) y \clave{argumentos} (las razones que la apoyan).
Puede tener también \clave{contraargumentos} (objeciones que el autor reconoce) y \clave{refutaciones} (sus respuestas).
Dibujarlo como un árbol aclara todo:

\figura{arbolarg}{El árbol de un argumento: la tesis arriba y sus razones debajo.}

\subsubsection{Tipos de argumento}
\begin{itemize}
  \item \textbf{De datos o ejemplos}: hechos, cifras, casos.
  \item \textbf{De autoridad}: lo que dice un experto o una institución.
  \item \textbf{De causa-consecuencia}: «si hacemos X, pasará Y».
  \item \textbf{De analogía}: «así como pasa en A, pasa en B».
  \item \textbf{De valores}: apela a lo justo, lo bueno, lo que todos estimamos.
\end{itemize}

\begin{regla}
\begin{itemize}
  \item La tesis responde a «¿qué quiere que yo piense?»; los argumentos, a «¿por qué debería pensarlo?».
  \item Prueba: la tesis es la que queda apoyada por las demás ideas, no la que apoya.
  \item Distingue la objeción que el autor presenta de su propia postura.
\end{itemize}
\end{regla}

\begin{texto}{Volver a cocinar}{Ensayo · texto original escrito para este libro}
\pa Cocinar en casa debería volver a ser parte de la educación de todos, hombres y mujeres. No se trata de nostalgia, sino
de salud, de economía y de autonomía.

\pa De salud, porque quien cocina sabe qué come: cuánta sal, cuánto azúcar, qué grasas. Los productos ultraprocesados,
que llenan las alacenas de quienes no cocinan, suelen esconder cantidades de esos ingredientes que nadie pondría a mano.

\pa De economía, porque un plato preparado en casa cuesta, en general, una fracción de lo que vale uno comprado. Una
familia que cocina sus almuerzos de la semana puede ahorrar una suma considerable a fin de mes.

\pa Se dirá que no hay tiempo. Es una objeción seria: las jornadas son largas y los trayectos también. Pero cocinar no
exige horas diarias; exige planear. Una tarde de domingo bien aprovechada puede resolver buena parte de la semana.

\pa Y de autonomía, finalmente, porque quien sabe cocinar depende menos de otros para algo tan básico como alimentarse.
Enseñarlo en el colegio sería enseñar una forma de libertad.
\end{texto}

\subsection{Preguntas comentadas}
\begin{comentada}{¿Qué función cumple el cuarto párrafo en la argumentación?}
\mal{A}{Presenta la tesis.}{La tesis está en el \S\,1.}
\mal{B}{Agrega un argumento de salud.}{Ese es el \S\,2.}
\bien{C}{Presenta una objeción (la falta de tiempo) y la responde.}{«Se dirá que…», «Es una objeción seria», «Pero…».}
\mal{D}{Concluye.}{La conclusión está en el \S\,5.}
\end{comentada}

\begin{reto}
«Las ciudades deberían sembrar más árboles: refrescan las calles, reducen el ruido y mejoran el ánimo de quienes caminan
bajo ellos.» ¿Cuál es la tesis?
\opcionesl{Los árboles refrescan las calles.}{Las ciudades deberían sembrar más árboles.}{Los árboles reducen el ruido.}{Caminar mejora el ánimo.}
\end{reto}

\begin{retoresuelto}
La tesis es lo que se defiende: «las ciudades deberían sembrar más árboles». Las otras tres opciones son las razones.
\textbf{Opción B.}
\end{retoresuelto}

\taller
\piso{Calienta}{el árbol}
\begin{ejercicios}
\item Dibuja el árbol del texto. \sol{Tesis (\S\,1): cocinar debe ser parte de la educación. Argumentos: salud (\S\,2), economía (\S\,3), autonomía (\S\,5). Objeción y refutación: falta de tiempo / planear (\S\,4).}
\end{ejercicios}
\piso{Como en el examen}{opción múltiple}
\begin{ejercicios}
\item El argumento del \S\,2 se apoya principalmente en
\opcionesl{la autoridad de un médico.}{el conocimiento de lo que se come al cocinar.}{una encuesta.}{una analogía con la escuela.}
\sol{B.}
\item ¿Cuál afirmación debilitaría más el argumento del \S\,3?
\opcionesl{Muchos productos ultraprocesados tienen mucha sal.}{En ciertos lugares, comprar comida preparada resulta más barato que cocinar.}{Cocinar es divertido.}{Los colegios tienen cocinas.}
\sol{B. Contradice la premisa económica.}
\item La palabra «libertad» del final conecta con el argumento de
\opcionesd{salud}{economía}{autonomía}{tiempo}
\sol{C.}
\end{ejercicios}

\bitacora

\leccion{pu-ag-2}{Intención y tono}{Juego de leer la voz detrás de las palabras.}

\tecnica{La técnica: la intención y la voz}
La \clave{intención} es lo que el autor quiere lograr: informar, explicar, persuadir, criticar, advertir, conmover,
divertir. El \clave{tono} es la actitud con que lo hace: objetivo, entusiasta, crítico, irónico, nostálgico, preocupado,
humorístico. Ambos se deducen de las palabras elegidas (adjetivos, adverbios, signos de exclamación), de lo que se destaca y
de lo que se calla.

\begin{center}\small
\begin{tabularx}{\linewidth}{@{}l>{\raggedright\arraybackslash}X@{}}
\toprule
\textbf{Tono} & \textbf{Pistas}\\\midrule
objetivo & datos, verbos neutros, ausencia de adjetivos valorativos\\
crítico & adjetivos negativos, preguntas que cuestionan\\
irónico & elogio aparente de algo malo; contraste entre lo dicho y la situación\\
nostálgico & pasado idealizado, «antes», «ya no»\\
preocupado & riesgos, advertencias, «si no hacemos algo»\\\bottomrule
\end{tabularx}
\end{center}

\figura{tono}{Cinco tonos frecuentes y la pista que los delata.}

\begin{regla}
\begin{itemize}
  \item Intención: ¿qué quiere que el lector haga, piense o sienta?
  \item Tono: subraya los adjetivos y adverbios valorativos.
  \item Elige la intensidad justa: «crítico» no es «furioso».
\end{itemize}
\end{regla}

\begin{texto}{El río que fue nuestro}{Columna · texto original escrito para este libro}
\pa Cuando yo era niño, los domingos bajábamos al río con una olla de sancocho y una pelota. Las mamás lavaban en las
piedras, los niños aprendíamos a nadar en los charcos y nadie pensaba que el agua pudiera acabarse.

\pa Hoy el río es una franja de piedras secas la mitad del año. Las curtiembres de arriba vierten lo que les sobra, y en
las orillas crecen bolsas plásticas en lugar de guaduales. Los niños del barrio lo conocen solo por las fotos de sus
abuelos.

\pa No escribo para lamentarme. Escribo porque todavía estamos a tiempo: hay ríos que se recuperaron cuando las
comunidades decidieron vigilarlos. Pero eso exige que dejemos de mirar hacia otro lado. El río no volverá solo.
\end{texto}

\subsection{Preguntas comentadas}
\begin{comentada}{¿Cuál es la intención principal del autor?}
\mal{A}{Contar anécdotas de su infancia para entretener.}{Las anécdotas están al servicio de otra cosa.}
\mal{B}{Informar sobre la química del agua.}{\emph{No lo dice}.}
\bien{C}{Advertir sobre el deterioro del río e invitar a actuar para recuperarlo.}{\S\,3: «todavía estamos a tiempo»,
«dejemos de mirar hacia otro lado».}
\mal{D}{Criticar a los abuelos.}{\emph{Dice otra cosa}.}
\end{comentada}

\begin{reto}
«Qué gran idea cerrar la única biblioteca del barrio justo cuando empiezan los exámenes.» El tono de la frase es
\opcionesd{entusiasta}{irónico}{neutral}{informativo}
\end{reto}

\begin{retoresuelto}
Se elogia («qué gran idea») algo claramente malo (cerrar la biblioteca en exámenes): es ironía. \textbf{Opción B.}
\end{retoresuelto}

\taller
\piso{Como en el examen}{opción múltiple}
\begin{ejercicios}
\item El tono del \S\,1 es predominantemente
\opcionesd{nostálgico}{irónico}{técnico}{indiferente}
\sol{A.}
\item La frase «No escribo para lamentarme» \pf{3} marca un cambio hacia un tono
\opcionesd{de llamado a la acción}{humorístico}{resignado}{neutral}
\sol{A.}
\item «En las orillas crecen bolsas plásticas en lugar de guaduales» \pf{2} expresa una crítica mediante
\opcionesl{un dato estadístico.}{una imagen que contrasta lo que debería haber con lo que hay.}{una cita de autoridad.}{una pregunta.}
\sol{B.}
\end{ejercicios}

\bitacora

\leccion{pu-ag-3}{Hechos y opiniones}{Juego de separar lo que se puede comprobar de lo que se valora.}

\tecnica{La técnica: ¿se puede comprobar?}
Un \clave{hecho} es una afirmación que se puede verificar (contando, midiendo, consultando un registro), sea verdadera o
falsa. Una \clave{opinión} expresa una valoración, un gusto o un juicio, y otra persona podría razonablemente pensar
distinto. Marcas de opinión: adjetivos valorativos (mejor, injusto, hermoso), verbos como «creo», «considero»,
expresiones como «debería», «es necesario».

\subsubsection{Mezclas}
Muchas oraciones mezclan ambas cosas: «La nueva vía, que costó 40\,000 millones, es un desperdicio» contiene un hecho (el
costo) y una opinión (el desperdicio). Las preguntas suelen pedir separar.

\figura{hechoopinion}{El termómetro da un hecho; lo que cada uno siente es opinión.}

\begin{regla}
\begin{itemize}
  \item ¿Se podría comprobar? Hecho. ¿Alguien podría valorarlo distinto? Opinión.
  \item Un hecho falso sigue siendo del tipo «hecho» (verificable).
  \item Las cifras y fechas suelen ser hechos; los adjetivos valorativos, opiniones.
\end{itemize}
\end{regla}

\begin{texto}{El nuevo estadio}{Editorial · texto original escrito para este libro (la ciudad y las cifras son imaginarias)}
\pa La ciudad inauguró ayer su nuevo estadio, con capacidad para 35\,000 personas y una inversión de 280\,000 millones de
pesos. La obra tardó cuatro años, uno más de lo previsto.

\pa Es, sin duda, un escenario hermoso. Pero cabe preguntarse si era la prioridad. Con ese dinero se habrían podido
construir varios colegios en los barrios que hoy tienen salones con cuarenta y cinco estudiantes.

\pa Los defensores del proyecto dicen que el estadio traerá turismo y empleo. Ojalá sea así. Por ahora, lo único seguro es
que la deuda la pagaremos todos durante los próximos quince años.
\end{texto}

\subsection{Preguntas comentadas}
\begin{comentada}{¿Cuál de las siguientes afirmaciones del texto es un hecho?}
\mal{A}{«Es, sin duda, un escenario hermoso.»}{Valoración estética: opinión.}
\mal{B}{«cabe preguntarse si era la prioridad.»}{Juicio del editorialista.}
\bien{C}{La obra tardó cuatro años, uno más de lo previsto.}{Se puede verificar con los registros de la obra.}
\mal{D}{«Ojalá sea así.»}{Deseo.}
\end{comentada}

\begin{reto}
¿Cuál de las siguientes es una opinión y no un hecho?
\opcionesl{El parque tiene doce bancas.}{El parque fue inaugurado en 2015.}{El parque es el lugar más agradable de la ciudad.}{El parque tiene un lago artificial.}
\end{reto}

\begin{retoresuelto}
«El lugar más agradable» es una valoración: otros podrían preferir otro lugar. Las demás se pueden verificar.
\textbf{Opción C.}
\end{retoresuelto}

\taller
\piso{Como en el examen}{opción múltiple}
\begin{ejercicios}
\item «El estadio traerá turismo y empleo» \pf{3} es
\opcionesl{un hecho comprobado.}{una predicción de los defensores, todavía no verificada.}{la opinión del editorialista.}{un dato de la obra.}
\sol{B.}
\item La postura del editorial frente al estadio es
\opcionesd{entusiasta}{crítica}{neutral}{indiferente}
\sol{B. Cuestiona la prioridad del gasto.}
\item ¿Qué hecho usa el editorial para apoyar su crítica?
\opcionesl{Que el estadio es hermoso.}{Que hay salones con cuarenta y cinco estudiantes.}{Que habrá turismo.}{Que a la gente le gusta el fútbol.}
\sol{B.}
\end{ejercicios}

\bitacora

\evaluacion
\begin{texto}{La jornada única}{Columna de opinión · texto original escrito para este libro}
\pa Ampliar la jornada escolar hasta la tarde es una buena idea, pero solo si cambia lo que se hace en ella. Más horas de
lo mismo no producen mejores estudiantes; producen estudiantes más cansados.

\pa Los defensores de la jornada única señalan, con razón, que mantiene a los jóvenes en un entorno protegido y que
permite enseñar artes, deportes y refuerzos que hoy no caben en el horario. Esos beneficios son reales.

\pa El problema es que, en muchos colegios, la jornada se amplió sin comedor, sin más docentes y sin espacios nuevos. El
resultado son tardes de «tiempo muerto» en las que los estudiantes esperan a que suene el timbre.

\pa Antes de celebrar las horas, deberíamos preguntar qué se enseña en ellas. Una hora bien usada vale más que tres horas
de espera.
\end{texto}
\begin{ejercicios}
\item La tesis del autor es que la jornada única
\opcionesl{debe eliminarse.}{es buena solo si cambia el contenido de esas horas.}{es perfecta tal como está.}{es mejor que cualquier otra medida.}
\sol{B. (4.1)}
\item El \S\,2 cumple la función de
\opcionesl{reconocer los argumentos de los defensores.}{refutar a los defensores.}{presentar datos de un estudio.}{concluir.}
\sol{A. (4.1)}
\item Las comillas en «tiempo muerto» \pf{3} y el tono del párrafo expresan
\opcionesd{una crítica}{entusiasmo}{humor sin intención}{neutralidad}
\sol{A. (4.2)}
\item ¿Cuál de estas afirmaciones es una opinión del autor?
\opcionesl{En muchos colegios la jornada se amplió sin comedor.}{Una hora bien usada vale más que tres horas de espera.}{La jornada única permite enseñar deportes.}{Los defensores señalan que es un entorno protegido.}
\sol{B. (4.3)}
\end{ejercicios}
